\subsection{\textbf{Select Papers}}
In this survey, we aim to provide a comprehensive review of research on Python
type analysis and type checking. To identify relevant studies, we conducted a
systematic search of the literature using keywords such as “type analysis,”
“static type checking,” and “dynamic type checking.” From the retrieved
studies, we selected only those that focus on Python or explicitly address the
dynamic characteristics of Python programs in their type inference or type
checking techniques.

Based on these criteria, the survey includes traditional static type
analyses~\cite{static2014}, constraint-based approaches~\cite{MaxSMT2018},
machine learning–driven methods~\cite{pyinfer2021,lambdanet2020,typeT52023},
and hybrid techniques combining static analysis and
learning~\cite{hityper2022,static2023,dlinfer2023}. 
In addition, we also consider studies that aim to predict or repair Python
runtime type errors~\cite{pyter2022,pyty2024,pyinder2024,quac2024}, covering
research that extends beyond pure type inference to practical program
improvement.

The selected studies collectively address diverse aspects of Python type
analysis, including handling dynamic features, compound data types, and flow
sensitivity. They represent a broad spectrum of methodologies—static, dynamic,
and learning-based approaches—and are thus well-suited for a comprehensive
comparison and discussion of the current state and future directions of Python
type checking research.
